\chapter{Base Tecnológica}
\label{chap:basetecnoloxica}

\lettrine{A}{} continuación se presenta la base tecnológica de este proyecto, manteniendo la división empleada en el trabajo previo \cite{TFGAdrianGarcia2025}. La diferencia principal es que, dentro de cada bloque, las tecnologías que ya existían se resumen de forma breve, mientras que las tecnologías nuevas o eliminadas se explican con más detalle para reflejar mejor la evolución de esta etapa.

\section{Lenguajes}

\subsection{Java}
Java \cite{java} sigue siendo el lenguaje principal del backend por su madurez, su integración natural con Spring y su adecuación al tipo de lógica que requiere la aplicación.

\subsection{SQL}
\acrfull{sql} \cite{sql} se mantiene como base para definir y consultar los datos persistidos.

\subsection{HTML y CSS}
\acrfull{html} y \acrfull{css} \cite{html_css} continúan siendo la base de la estructura y presentación de la interfaz web.

\subsection{JavaScript}
\acrfull{javascript} se conserva en la capa frontend por su papel en la lógica de interacción y en la construcción de componentes de React \cite{javascript}.

\subsection{JSX}
\acrfull{jsx} \cite{jsx} sigue siendo el formato de escritura de los componentes de React, ya que permite combinar la estructura visual con la lógica de la interfaz de forma clara.

\subsection{LaTeX}
\acrfull{latex} \cite{latex} se mantiene para la redacción y composición de la memoria.

\subsection{JSON}
\acrfull{json} \cite{json} continúa siendo el formato de intercambio de datos entre cliente y servidor.

\section{Frameworks y librerías servidor}

\subsection{Spring}
\acrfull{spring} \cite{spring} permanece como núcleo del backend por su robustez y por facilitar el acceso a datos y la lógica de negocio.

\subsection{Spring Boot}
\acrfull{springboot} \cite{springboot} sigue simplificando la configuración del backend y el despliegue de la aplicación.

\subsection{Spring Data}
\acrfull{springdata} \cite{springdata} se mantiene como capa de abstracción sobre el acceso a datos, reduciendo la necesidad de consultas manuales.

\subsection{JPA}
\acrfull{jpa} \cite{jpa} continúa modelando la persistencia de las entidades del dominio.

\subsection{Hibernate}
Hibernate \cite{hibernate} sigue siendo la implementación de referencia para el trabajo con entidades, relaciones y carga perezosa.

\subsection{JDBC}
\acrfull{jdbc} \cite{jdbc} se conserva como soporte de bajo nivel cuando se necesita acceso directo a la base de datos.

\subsection{OpenAPI}
OpenAPI \cite{OpenAPI} es una de las incorporaciones más relevantes de este trabajo. Se utiliza para definir primero el contrato de la API en YAML y generar a partir de él controladores y DTOs. Este enfoque \textit{API-first} permite que la documentación, el backend, el frontend y Android compartan una misma definición de los servicios, reduciendo desajustes entre implementación y consumo.

\section{Frameworks y librerías web}

\subsection{React}
\acrfull{react} \cite{react} se conserva como base de la interfaz web por su capacidad para construir una aplicación modular y mantenible.

\subsection{Redux}
\acrfull{redux} \cite{redux} sigue siendo el contenedor central del estado global de la aplicación web.

\subsection{RTK Query}
\acrfull{rtkquery} \cite{rtkquery} se mantiene como solución para el acceso a la API, al simplificar caché, consultas y mutaciones.

\subsection{Material UI}
Material UI \cite{materialui} continúa proporcionando los componentes visuales principales de la aplicación web.

\subsection{React Map GL}
React Map GL \cite{reactmapgl} se elimina como dependencia explícita porque la integración cartográfica se centraliza sobre MapLibre \cite{maplibre}. Esta sustitución reduce dependencia de una capa anterior y unifica el tratamiento del mapa sobre una tecnología más abierta y alineada con el resto del proyecto.

\subsection{Open Weather Map}
Open Weather Map \cite{openweathermap} se mantiene como fuente externa de información meteorológica útil en el contexto de emergencias.

\subsection{Mapbox}
Mapbox \cite{mapbox} se descarta como solución principal para reducir la dependencia de claves API y de planes de prueba o gratuitos. En este trabajo, la visualización geográfica se apoya en MapLibre, que se adapta mejor al enfoque abierto del proyecto.

\section{Pruebas}

\subsection{JUnit}
\acrfull{junit} \cite{junit} se conserva como herramienta principal de pruebas unitarias.

\section{Protocolos}

\subsection{HTTP/S}
\acrfull{http} y \acrfull{https} \cite{http} continúan siendo la base de la comunicación cliente-servidor.

\subsection{REST}
\acrfull{rest} \cite{rest} sigue siendo el estilo principal para exponer recursos y operaciones entre cliente y servidor.

\section{Herramientas de desarrollo}

\subsection{IntelliJ}
\acrfull{ide} IntelliJ \cite{intellij} sigue siendo la herramienta principal de desarrollo del backend.

\subsection{Postman}
Postman \cite{postman} se mantiene como herramienta de pruebas y exploración de endpoints.

\subsection{VSCode}
VSCode \cite{vscode} se conserva como editor de apoyo para la parte web y la documentación.

\subsection{DBeaver}
DBeaver \cite{dbeaver} continúa siendo la herramienta de consulta y administración de la base de datos.

\subsection{Git}
Git \cite{git} se mantiene como sistema de control de versiones del proyecto.

\subsection{Apache Maven}
Apache Maven \cite{maven} sigue siendo la base del ciclo de construcción del backend y también soporta la generación de OpenAPI.

\subsection{Create React App}
Create React App \cite{create-react-app} permanece como base de arranque del proyecto web.

\section{Sistemas de gestión de base de datos}

\subsection{H2}
H2 \cite{h2} se elimina como base de datos principal porque en este trabajo se trabaja directamente sobre PostgreSQL y PostGIS para acercar el entorno de desarrollo al de ejecución real. En el trabajo previo era útil para pruebas rápidas y desarrollo local, pero ahora resulta más coherente usar la misma base de datos que la aplicación final.

\subsection{PostgreSQL}
PostgreSQL \cite{postgresql} se mantiene como sistema de gestión principal por su estabilidad y su integración con el resto de componentes.

\subsection{PostGIS}
PostGIS \cite{postgis} sigue siendo esencial para soportar las operaciones geoespaciales del proyecto.

\section{Tecnologías de alojamiento de servicios}

\subsection{Tomcat}
Tomcat \cite{tomcat} se mantiene como contenedor de ejecución del backend.

\subsection{Webpack}
Webpack \cite{webpack} continúa formando parte de la cadena de construcción de la interfaz web.

\subsection{Serve}
Serve \cite{serve} se conserva como herramienta sencilla para servir la versión compilada del frontend durante pruebas o despliegues locales.

\section{Tecnologías nuevas}

\subsection{Kotlin}
Kotlin \cite{kotlin} se incorpora como lenguaje principal de la aplicación Android por su concisión, su interoperabilidad con Java y su buena integración con el ecosistema Android.

\subsection{Gradle}
Gradle \cite{gradle} se añade como sistema de construcción habitual en Android, ya que facilita la gestión de módulos, dependencias y variantes de compilación.

\subsection{Jetpack Compose}
\acrfull{compose} \cite{jetpack_compose} se incorpora para crear interfaces Android de forma declarativa y más mantenible, adaptándose mejor a una aplicación con múltiples pantallas y navegación compleja.

\subsection{Material Design}
\acrfull{mat} \cite{material_design} se incorpora como guía visual para la app móvil, con el objetivo de mantener una experiencia de usuario coherente y reconocible.

\subsection{HttpClient}
HttpClient \cite{httpclient} se añade como biblioteca para realizar peticiones HTTP desde Android y centralizar el acceso a los servicios del backend.

\subsection{Firebase Cloud Messaging}
\acrfull{fcm} \cite{firebase} se incorpora para el envío de notificaciones push, que es una funcionalidad clave en el contexto móvil de la plataforma.

\subsection{Firebase Messaging Service}
\acrfull{fms} \cite{firebase_messaging_service} se añade para gestionar la recepción de mensajes procedentes de \acrshort{fcm} y reaccionar a notificaciones o eventos remotos.

\subsection{Testcontainers}
\acrfull{testcontainers} \cite{testcontainers} se incorpora para ejecutar pruebas de integración sobre una base de datos aislada y reproducible. Esta tecnología está basada en Docker \cite{docker} y permite levantar contenedores específicos para cada prueba o conjunto de pruebas, garantizando un entorno controlado y evitando interferencias entre pruebas.
